\section{Introduction}
\label{sec:intro}

In 2023 the Prime Video engineering team reported that moving an audio/video monitoring service from a
distributed design built on serverless components into a single process ``reduced our infrastructure cost by
over 90\%'' and increased its scaling capabilities~\cite{Kolny2023}. Su et al.\ analyse this case as one of four
documented moves from microservices back to a monolith, with cost, complexity, scalability, performance and
organisation as the recurring reasons~\cite{Su2024}. Such reports do not settle whether decomposition is a
good idea. They do show that the performance side of the decision is still argued from individual cases.

The published comparisons do not settle it either. Reviews of the literature on monolith-to-microservice
migration find conflicting results~\cite{Rodrigues2024} and a lack of comparative studies of architectural
styles~\cite{RodriguesJSS2025}. Much of the disagreement can be traced to method. Most comparisons report
throughput or mean response time rather than the latency distribution. Many use closed-loop load generators,
which slow down with the system they measure and so hide the tail~\cite{Schroeder2006,Tene2015}. None that we know
of splits end-to-end latency into the layers where distribution adds cost. And the resources given to each
architecture are rarely held equal or even reported, although a microservice deployment of the same code runs
more processes, each with its own runtime, connection pools and share of the CPU.

This article reports a controlled comparison designed around those four problems. It implements and executes a
measurement protocol we presented earlier\footnote{As a talk at a university conference at the Kazakh-British
Technical University, 2026, without published proceedings. \textbf{[TODO(author): conference name and date.]}},
and it revises that protocol where building it showed parts of it to be unsound. The testbed runs one e-commerce application in four topologies: a monolith (A), four
coarse-grained services (B), eight fine-grained services (C), and the same eight services with an Istio service
mesh (D). All four run the same binary with identical business code; only the mapping of modules to processes
differs. Every configuration receives the same CPU budget. Each process reports, in a W3C \texttt{Server-Timing}
header propagated through the call tree, how long it spent in application code, database calls, serialisation,
framework dispatch and waiting for its callees, which yields an exact per-request decomposition of end-to-end
latency without synchronised clocks. Latency is measured under open-loop load, capacity under closed-loop load,
and the Universal Scalability Law (USL)~\cite{Gunther2007} is fitted to the latter.

The study answers four research questions:
\begin{description}
\item[RQ1] How much tail latency does decomposition add at low load, when every configuration has spare
capacity?
\item[RQ2] Under an equal resource budget, does any decomposed configuration reach a higher throughput than the
monolith, and where do their capacity curves cross? Does the answer change when the budget is allocated in
proportion to demand, or when an autoscaler can add replicas up to a shared ceiling?
\item[RQ3] In which layers does the additional latency of decomposition arise, and what limits each
configuration's capacity?
\item[RQ4] How much does a sidecar service mesh add, and does its cost grow with the depth of the call graph?
\end{description}
RQ1, RQ2 and RQ4 correspond to hypotheses H1--H3 of the protocol, each stated with an explicit falsification
condition (Section~\ref{sec:hypotheses}).

In \the\numexpr\xo{cells}+\xo{prop/cells}+\xo{work/cells}+\xo{workprop/cells}+\xo{eks/cells}+\xo{hpa/cells}\relax{} fresh deployments the answers were the same in both environments. Decomposition raised
low-load p99 latency, mostly through inter-service transit rather than serialisation (H1 supported). Under an equal
budget no decomposed configuration overtook the monolith at any load (H2 falsified). Each was limited by the CPU
quota of its edge service, whose throttling made tail latency erratic before its average utilisation reached the
limit, while most of the budget stayed unused. Allocating the budget by demand and giving every operation more work
narrowed the gap but did not close it.
Autoscaling under a shared CPU ceiling did not change this: the monolith sustained several times the load of the
decomposed configurations on half the CPU they requested at the ceiling.
The Istio mesh added a fixed cost per call, \xo{H3/share_deep}\% (local) and \xo{eks/H3/share_deep}\% (cloud) of
latency on a six-hop path, and used more CPU than the application (H3 supported).


The contributions are:
\begin{enumerate}
\item a testbed design in which configurations differ only in deployment topology, run under an equal CPU budget,
and are checked for identical responses on every deployment;
\item a per-layer latency accounting based on \texttt{Server-Timing} that assigns every millisecond of end-to-end
latency to one layer and needs no clock synchronisation, together with cgroup-level CPU and throttling data that
explain the capacity limits;
\item measurements from \xo{cells} fresh deployments on a local cluster and \xo{eks/cells} on Amazon EKS,
analysed with the repetition as the unit of replication. They test H1--H3 and locate (or rule out) the crossover
under an equal budget, under demand-proportional allocation and under autoscaling;
\item the testbed, campaign runner, raw data, analysis pipeline and an interactive dashboard, from which every
number in this article is generated.
\end{enumerate}

Section~\ref{sec:background} explains the mechanisms by which distribution costs time and the models used to
analyse them. Section~\ref{sec:related} reviews related empirical work. Sections~\ref{sec:testbed}
and~\ref{sec:method} describe the testbed and the experimental method. Section~\ref{sec:results} reports the
results, Section~\ref{sec:discussion} discusses what they mean for practitioners, and
Section~\ref{sec:threats} discusses threats to validity.
